\documentclass[11pt]{article}
\usepackage[margin=1in]{geometry}
\usepackage{amsmath, amssymb, amsthm}
\usepackage{mathtools}
\usepackage{hyperref}

\newtheorem{theorem}{Theorem}
\newtheorem{lemma}[theorem]{Lemma}
\newtheorem{proposition}[theorem]{Proposition}
\newtheorem{corollary}[theorem]{Corollary}
\newtheorem{conjecture}[theorem]{Conjecture}
\theoremstyle{definition}
\newtheorem{definition}[theorem]{Definition}
\newtheorem{remark}[theorem]{Remark}

\newcommand{\Sym}{\mathfrak{S}}
\newcommand{\Hq}{\mathcal{H}_q}
\newcommand{\la}{\lambda}
\newcommand{\hla}{\hat{\lambda}}
\newcommand{\Om}{\Omega}
\newcommand{\Omstar}{\Omega^{*}}
\newcommand{\vT}{v_T}
\newcommand{\SYT}{\mathrm{SYT}}

\title{Residual diagonals for $\la = (\hla, 1^q)$:\\
refutation of the Kazhdan--Lusztig cell hypothesis\\
and an empirical column-$1$ SC discriminator\\
at $\hla = (3, 3)$}
\author{Clio}
\date{2026-05-19}

\begin{document}
\maketitle

\begin{abstract}
The SC-at-$1$ trivial reduction closes $378/514$ of the empirically
tested converse Diagnostic 4 cases. The remaining open core is the
\emph{SR-at-$1$ residual}: zero diagonals $p_T(q) = 0$ for which the
only column-strict pairs in $T$ occur at indices $i \ge 2$. We make
three contributions. First, we test and refute the
Dai--Zhang-inspired hypothesis that the SR-at-$1$ residual is a union
of cells in the Hoefsmit/Murphy $W$-graph or the descent partition;
on $\la = (3,3,1,1,1)$, two descent classes each contain one residual
and one nonresidual member. Second, we propose and verify (on $50/50$
SR-at-$1$ tableaux of $(3,3,1,1,1)$ with $0$ errors) an explicit
combinatorial discriminator built from \emph{column-$1$} SC pairs at
the arithmetic-progression indices $\{3, 5, 7\}$. Third, we identify
a sign-cancellation phenomenon: three nonresidual tableaux with
``double SC at $\gamma$-position'' (entries $7, 8, 9$ all stacked in
column $1$) escape the discriminator's $\gamma$-clause, and the
non-continuation clause is essential. We retract the earlier
$\Om^*$-only count of $12$ for this shape: all $35$ residual zero
diagonals at $(3,3,1,1,1)$ are killed by $\Om$ \emph{and} $\Om^*$
simultaneously.
\end{abstract}

\section{Background and the residual problem}

\subsection{Setting and notation}
Fix $\la = (\hla, 1^q)$ with $\ell(\hla) \ge 2$, every row
$\hla_r \ge 2$, and $\tau(\hla) := \max(0, q + 2\ell(\hla) - 1 -
|\hla|) = 0$. Let $n = |\la|$, $\ell = \ell(\la)$. Let $V^\la$ be the
irreducible $\Hq(\Sym_n)$-module in the Hoefsmit (seminormal) basis
$\{v_T : T \in \SYT(\la)\}$; the Hecke generators $T_i$ satisfy
$(T_i - q)(T_i + 1) = 0$ and we set $S_i := T_i + 1$. Define the
chain factors and the diagnostic operator
\begin{align*}
R'_k &:= S_{k-1} S_{k-2} \cdots S_2 S_1, \\
L_k  &:= S_1 S_2 \cdots S_{k-2} S_{k-1}, \\
\Om   &:= R'_{\ell+1} R'_{\ell+2} \cdots R'_n, \\
\Omstar &:= L_n L_{n-1} \cdots L_{\ell+1}.
\end{align*}
The diagonal probe is
$p_T(q) := \langle v_T,\, \Om\, v_T\rangle \in \mathbb{Q}(q)$ with
respect to the Hoefsmit pairing.

\begin{definition}
For $T \in \SYT(\la)$, an \emph{SC pair} at index $i$ is a pair
$(i, i+1)$ with the cell containing $i+1$ directly below the cell
containing $i$ (same column). An \emph{SR pair} is the row analogue.
We say $T$ is \emph{SC-at-$1$} if the cell containing $2$ is $(2,1)$,
otherwise (since $2$ lies at $(1,2)$ or $(2,1)$ in any SYT) $T$ is
\emph{SR-at-$1$}.
\end{definition}

\subsection{The two facts in hand}
Two structural results frame the residual question.

\begin{itemize}
\item \textbf{Diagnostic 4 dichotomy} (forward direction proved
$2026$-$05$-$18$).\
$p_T(q) = 0 \;\Leftrightarrow\; \Om v_T = 0 \;\text{or}\;
\Omstar v_T = 0$. The forward direction is proved by the
two-sided sign-kill identity; the converse is open.
\item \textbf{Universal SC necessity} (proved $2026$-$05$-$18$).\
If $p_T(q) = 0$ then $T$ has at least one SC pair (Hoefsmit
positivity rigidity on the diagonal path). The converse fails: $T$
may have an SC pair with $p_T \ne 0$.
\end{itemize}

\subsection{The SC-at-$1$ trivial reduction}
Yesterday's note \texttt{2026-05-19-sc-at-1-trivial-reduction.tex}
showed that if $T$ is SC-at-$1$ then $S_1 v_T = 0$, hence
$R'_n v_T = 0$ (since $S_1$ is the rightmost factor of $R'_n$),
hence $\Om v_T = 0$. This closes $378/514 \approx 73.5\%$ of the
converse cases in our $12$-shape corpus.

\subsection{The residual: SR-at-$1$}
The remaining open problem isolates the $136$ zero diagonals
$T$ with $T$ SR-at-$1$, i.e.\ entries $1, 2$ both in row $1$. The
algebraic obstacle is that for SR-at-$1$ tableaux, $S_1$ does not
annihilate $v_T$, so the one-line argument fails. The full converse
of Diagnostic 4 reduces to:

\begin{quote}
For every $T \in \SYT(\la)$ with $T$ SR-at-$1$ and $p_T(q) = 0$,
either $\Om v_T = 0$ or $\Omstar v_T = 0$.
\end{quote}

This is the SR-at-$1$ residual problem. At $\la = (3,3,1,1,1)$, the
residual contains $35$ tableaux among $50$ SR-at-$1$ tableaux (the
other $15$ are SR-nonzero).

\section{The cell hypothesis}

\subsection{Motivation}
Dai--Zhang \cite{DaiZhang} construct \emph{recursive Gelfand
$W$-graphs} for the symmetric group, refining classical Kazhdan--
Lusztig cells. In the Hoefsmit/Murphy basis of $V^\la$, several
natural partitions on $\SYT(\la)$ arise:
\begin{itemize}
\item the Murphy graph (SCC partition on the basis transition graph),
\item the \emph{descent partition}: two SYTs are equivalent iff they
have the same descent set $D(T) := \{i : i+1 \text{ is strictly
below } i \text{ in } T\}$,
\item the dual-equivalence (Assaf) partition,
\item the (recursive) Gelfand $W$-graph cells of \cite{DaiZhang}.
\end{itemize}
The hopeful version of the SR-residual hypothesis runs:

\begin{conjecture}[Cell hypothesis]\label{conj:cell}
There exists a partition $\Pi$ of $\SYT(\la)$ arising from a
Kazhdan--Lusztig/Murphy-type $W$-graph such that the SR-at-$1$
residual is a union of blocks of $\Pi$.
\end{conjecture}

If true, Conjecture~\ref{conj:cell} would yield a uniform proof of
the SR-at-$1$ converse by cell-induction \cite{DaiZhang}. The
\emph{descent partition} is the finest natural partition compatible
with KL left cells, and the Murphy/Knuth dual-equivalence partition
within a fixed shape is one giant class (Haiman). So the descent
partition is the strongest test of Conjecture~\ref{conj:cell}.

\section{Refutation at $\la = (3,3,1,1,1)$}

\subsection{Setup}
We take $\la = (3, 3, 1, 1, 1)$, so $\hla = (3, 3)$, $q = 3$,
$\ell(\hla) = 2$, $|\hla| = 6$, $\tau = 0$, $n = 9$, and
\[
\Om = R'_3\, R'_4\, R'_5\, R'_6\, R'_7\, R'_8\, R'_9.
\]
The Specht module $V^\la$ has $N = 120$ basis vectors. Of these,
$70$ are SC-at-$1$ (each closed by the trivial reduction) and $50$
are SR-at-$1$. Among the SR-at-$1$ tableaux, $35$ are residual
($p_T = 0$) and $15$ are nonresidual ($p_T \ne 0$). All counts and
diagonal probes use the symbolic Hoefsmit basis specialised at
$q \in \{2, 3, 5, 7\}$; agreement at four distinct positive rationals
fixes the diagonal entry to $0$ in $\mathbb{Q}(q)$ at the degrees
encountered here.

\subsection{The decisive descent blocks}
The descent partition on $\SYT(3,3,1,1,1)$ has $95$ blocks ($21$
non-singleton). Two non-singleton descent blocks each contain one
SR-residual and one SR-nonresidual member:

\medskip
\noindent\textbf{Block } $D = (2, 4, 5, 7, 8)$ \textbf{ with members}
$\{T_{11}, T_{27}\}$\textbf{:}
\begin{itemize}
\item $T_{11} = (\,1,2,4 \mid 3,5,7 \mid 6 \mid 8 \mid 9\,)$. SC list
$= [\,8\,]$. SR list $= [\,1\,]$. $p_{T_{11}} \ne 0$.
\item $T_{27} = (\,1,2,7 \mid 3,4,8 \mid 5 \mid 6 \mid 9\,)$. SC list
$= [\,5, 7\,]$. SR list $= [\,1, 3\,]$. $p_{T_{27}} = 0$.
\end{itemize}

\noindent\textbf{Block } $D = (2, 4, 5, 6, 8)$ \textbf{ with members}
$\{T_{12}, T_{29}\}$\textbf{:}
\begin{itemize}
\item $T_{12} = (\,1,2,4 \mid 3,5,8 \mid 6 \mid 7 \mid 9\,)$. SC list
$= [\,\,]$. SR list $= [\,1\,]$. $p_{T_{12}} \ne 0$.
\item $T_{29} = (\,1,2,8 \mid 3,4,9 \mid 5 \mid 6 \mid 7\,)$. SC list
$= [\,5, 6, 8\,]$. SR list $= [\,1, 3\,]$. $p_{T_{29}} = 0$.
\end{itemize}

\medskip
In each block, both members are SR-at-$1$ (entries $1, 2$ in row $1$)
and share the same descent set, yet exactly one is residual. The
nonresidual member ($T_{11}$ or $T_{12}$) has no SC pair below index
$8$; the residual member has SC pairs at indices $\le 6$.

\subsection{Auxiliary checks}
The Murphy SCC partition on $\SYT(3,3,1,1,1)$ in the Hoefsmit basis
is the trivial giant component (one block). Thus the Murphy partition
contains the residual set vacuously but provides no separating power.
Among partitions strictly coarser than the trivial one, the descent
partition is the finest reasonable cell candidate, and it already
fails.

\begin{theorem}[Refutation of the cell hypothesis at $(3,3,1,1,1)$]
\label{thm:cell-refuted}
The SR-at-$1$ residual on $\la = (3,3,1,1,1)$ is not a union of
descent classes. In particular, it is not a union of blocks of any
partition coarser than or equal to the descent partition; this rules
out the Murphy and dual-equivalence partitions on this shape as
candidates for Conjecture~\ref{conj:cell}.
\end{theorem}

\begin{proof}
The blocks $\{T_{11}, T_{27}\}$ and $\{T_{12}, T_{29}\}$ each
contain exactly one residual member. A partition cannot have the
residual set as a union of blocks if some block splits between
residual and nonresidual.
\end{proof}

\begin{remark}
Theorem~\ref{thm:cell-refuted} does not rule out finer
KL-recursive partitions (e.g.\ Gelfand $W$-graph cells from
\cite{DaiZhang} or asymptotic two-sided cells). It does say that the
naive partitions accessible from the Hoefsmit basis alone do not
encode the residual stratum, and any uniform cell-theoretic proof
must descend below the descent partition.
\end{remark}

\section{The column-$1$ SC discriminator at $\hla = (3, 3)$}

\subsection{Statement}
For each $T \in \SYT(3,3,1,1,1)$, write $T(r, c)$ for the entry in
row $r$, column $c$ (so column $1$ has entries
$T(1,1) = 1, T(2,1), T(3,1), T(4,1), T(5,1) = 9$, with the last
forced because $9$ is the largest entry and $(5, 1)$ is the only
extremal cell with row index $5$).

\begin{theorem}[Empirical, $\la = (3,3,1,1,1)$]\label{thm:emp}
For every $T \in \SYT(3,3,1,1,1)$ with entries $1, 2$ in row $1$
(i.e.\ $T$ SR-at-$1$), $p_T(q) = 0$ if and only if at least one of
the following holds:
\begin{itemize}
\item[\textup{($\alpha$)}] entries $3, 4$ both in column $1$,
equivalently $T(3, 1) = 4$;
\item[\textup{($\beta$)}] entries $3, 5, 6$ in column $1$,
equivalently $T(2, 1) = 3$, $T(3, 1) = 5$, $T(4, 1) = 6$;
\item[\textup{($\gamma$)}] entries $7, 8$ in column $1$ but $9$ not
in column $1$ above the bottom, equivalently $T(4, 1) = 7,
T(5, 1) = 8$ with the constraint forcing $9 \in (2, 3)$ (so
$T(5,1)=8$ is below $T(4,1)=7$, and $9$ sits in row $2$ not in
column $1$).
\end{itemize}
This statement was verified on all $50$ SR-at-$1$ tableaux with $0$
counterexamples.
\end{theorem}

\subsection{Arithmetic structure}
The discriminating SC-pair indices are
\[
\{3, 5, 7\} = \{q,\; q + 2,\; q + 4\}, \qquad q = 3.
\]
This is the column-$1$ promotion arithmetic of the SRD chain
$\Om = R'_3 \cdots R'_9$, with column-$1$ transitions spaced by $2$.
The three clauses ($\alpha$), ($\beta$), ($\gamma$) correspond to
which factor of $\Om$, in the right-to-left order
$R'_9, R'_8, \ldots, R'_3$, first annihilates the running vector
$\Om_{\ge k} v_T$:
\begin{center}
\begin{tabular}{c|c|c|c}
clause & col-$1$ SC pair & first-kill $R'_k$ & count \\
\hline
$\alpha$ & $(3, 4)$ at indices $(2,1), (3,1)$ & $R'_8$ & $20$ \\
$\beta$ & $(5, 6)$ at indices $(3,1), (4,1)$ & $R'_7$ & $9$ \\
$\gamma$ & $(7, 8)$ at indices $(4,1), (5,1)$, $9 \notin $ col $1$ & $R'_6$ & $6$ \\
\end{tabular}
\end{center}
The three counts $20 + 9 + 6 = 35$ match the residual cardinality
exactly.

\subsection{Refuted simpler conjectures}
Several simpler discriminators were tested against the same data and
refuted.

\medskip
\noindent\textbf{(Conj A) SC index threshold.} Try
$p_T = 0 \Leftrightarrow \min(\mathrm{SC}(T)) \le t$ or
$\max(\mathrm{SC}(T)) \le t$ for some uniform $t$. Refuted: on the
$50$ SR-at-$1$ tableaux, the zero-set values of $\min \mathrm{SC}$ are
$\{3, 4, 5, 6, 7\}$ and the nonzero-set values are
$\{4, 5, 6, 7, 8\}$; the overlap kills any monotone-threshold rule in
either direction. The same overlap occurs with $\max \mathrm{SC}$.

\medskip
\noindent\textbf{(Conj B) SC pair with prefix $\{1, \ldots, i\}$
filling row $1$.} Try $p_T = 0 \Leftrightarrow$ $T$ has an SC pair at
some index $i$ such that $\{1, 2, \ldots, i\}$ fills row $1$ from the
left. Vacuous on SR-at-$1$: row $1$ has length $3$, so the prefix
condition demands $i \le 3$; SR-at-$1$ forbids SC at $i = 1$; SC at
$i = 2$ requires $2$ directly above $3$ in column $2$, impossible
when row $1$ starts $1, 2, X$ with $X \ge 3$. The conjecture predicts
$p_T = 0$ on no SYT and is therefore false on all $35$ residual
tableaux.

\medskip
\noindent\textbf{(Conj C) Row-$1$ small-entry condition.} Try
$p_T = 0 \Leftrightarrow$ either row $1$ of $T$ equals $1, 2, 3$, or
row $1$ starts $1, 2$ and row $2$ starts $3, 4$ with cells $(2,1),
(2,2)$. Refuted from both sides: among the $35$ residual tableaux,
$25$ satisfy \emph{neither} sub-clause of (Conj C), and among the
$10$ tableaux satisfying one of the sub-clauses, $5$ are nonresidual.
This conjecture, originally suggested by the
[[omega-survivor-pattern]] memory note, captures only a small subset
of the residual.

\subsection{Verification details}
The discriminator was checked tableau-by-tableau, recording for each
$T$ the column-$1$ entries $T(2,1), T(3,1), T(4,1), T(5,1)$, the SC
list, the descent set, and the first-vanish stage of $\Om v_T$
under the right-to-left factorisation $\Om = R'_3 \cdots R'_9$. The
35 residual tableaux partition cleanly into the three clauses with
no overlap (no tableau satisfies two of $(\alpha, \beta, \gamma)$
simultaneously in the verification data), and the $15$ nonresidual
SR-at-$1$ tableaux satisfy none of $(\alpha, \beta, \gamma)$ except
for the three sign-cancellation cases discussed in
Section~\ref{sec:sign-cancel}.

The verification script and full table are at
\texttt{2026-05-19-V33111-residual-diagnostic.py} and
\texttt{2026-05-19-V33111-residual-diagnostic.md}.

\section{The sign-cancellation phenomenon}
\label{sec:sign-cancel}

The three SYTs $T_0, T_{10}, T_{20}$ have entries $7, 8, 9$ all
stacked in column $1$ (rows $3, 4, 5$). They therefore satisfy
\emph{both} the SC pair $(7, 8)$ in column $1$ and the SC pair
$(8, 9)$ in column $1$. The naive $\gamma$-clause without the
exclusion ``$9$ not in column $1$'' would predict $p_T = 0$ on these
three; in fact $p_T \ne 0$ for all three.

\begin{itemize}
\item $T_0 = (1,2,3 \mid 4,5,6 \mid 7 \mid 8 \mid 9)$: SC $= [7, 8]$,
$p_{T_0} \ne 0$.
\item $T_{10} = (1,2,4 \mid 3,5,6 \mid 7 \mid 8 \mid 9)$: SC
$= [7, 8]$, $p_{T_{10}} \ne 0$.
\item $T_{20} = (1,2,5 \mid 3,4,6 \mid 7 \mid 8 \mid 9)$: SC
$= [5, 7, 8]$, $p_{T_{20}} \ne 0$.
\end{itemize}

Interpretation: the column-$1$ SC pair $(7, 8)$ contributes a
$-1$-eigenvalue absorption in $R'_8$, which on its own would force
$R'_6 \cdots R'_8 v_T \in \ker S_7$. But the further SC pair
$(8, 9)$ in column $1$ contributes an additional $-1$-absorption in
$R'_9$, and the two interfere so that the net partial product
$R'_6 \cdots R'_9 v_T$ no longer lies in $\ker S_7$. The
``non-continuation'' clause in ($\gamma$) is therefore essential:
it forbids the SC pair $(8, 9)$ stacked immediately below $(7, 8)$
in column $1$, breaking the cancellation.

The analogous neutralisation does not occur for $\alpha$ or $\beta$
because the candidate ``next'' SC pairs $(4, 5)$ in column $1$ (for
$\alpha$) and $(6, 7)$ in column $1$ (for $\beta$) are never realised
in any SR-at-$1$ tableau of this shape: the SR-at-$1$ condition
together with the column-$1$ length-$5$ constraint forces $T(3, 1)
\in \{4, 5, 6, 7, 8\}$ varying, and $T(4, 1) > T(3, 1)$ with $T(5, 1)
= 9$. In particular $T(3, 1) = 4 \Rightarrow T(4, 1) \in
\{5, 6, 7, 8\}$, so the pair $(4, 5)$ in column $1$ is realised only
when $T(4, 1) = 5$ alongside $T(3, 1) = 4$, but then the pair $(5,
6)$ in column $1$ would require $T(5, 1) = 6 \ne 9$, contradiction.
A similar contradiction rules out $(6, 7)$ stacked immediately below
$(5, 6)$ in column $1$. The asymmetry between ($\gamma$) and
($\alpha, \beta$) is therefore an artefact of $9$ being the largest
entry and naturally landing at the bottom of column $1$.

\section{Open questions}

\subsection{Generalisation to other $\hla$}
Theorem~\ref{thm:emp} is verified on one shape. The arithmetic
structure ``column-$1$ SC pairs at indices in $\{q, q+2, q+4, \ldots\}$''
suggests a uniform rule:

\begin{conjecture}[Column-$1$ SC discriminator, general form]
\label{conj:general}
For $\la = (\hla, 1^q)$ with $\hla$ admitting an SR-at-$1$ stratum
and $\tau(\hla) = 0$, an SR-at-$1$ tableau $T$ satisfies $p_T = 0$ if
and only if $T$ has an SC pair in column $1$ at some index $i \in
\{q, q + 2, \ldots\}$ within the column-$1$ length range, subject to a
non-continuation condition at the bottom of column $1$.
\end{conjecture}

\noindent The smallest accessible test is $\la = (3, 3, 2)$, where
the prior probe found $5$ residual tableaux all satisfying SC at
index $3$ (the column-$1$ start for $\hla = (3, 3, 1)$); this is
consistent with Conjecture~\ref{conj:general} but vacuous in the
discriminator's nontrivial direction. The next informative test is
$\la = (4, 3, 1, 1, 1)$ or $\la = (3, 3, 1, 1, 1, 1)$, where the
column-$1$ promotion arithmetic predicts spacing $2$ but at shifted
base.

\subsection{Structural proof}
The discriminator should follow from a Hecke-algebraic identity of
the form: \emph{the partial chain product
$R'_{i+1} R'_{i+2} \cdots R'_n v_T$ lies in $\ker S_i$ whenever $T$
has an SC pair in column $1$ at index $i \in \{q, q + 2, q + 4\}$
without continuation}.
For ($\alpha$) with $i = 3$, this asks that
$R'_4 \cdots R'_9 v_T \in \ker S_3 = \mathrm{span}\{v_S : S \text{ has
SC at } 3\}$ when $T(2, 1) = 3, T(3, 1) = 4$. The natural route is to
analyse how $R'_k$ acts on the position of the letter $i + 1$ in the
support of the running vector (a ``letter-shift'' argument, by
analogy with the SC-at-$1$ trivial proof where $R'_n$ has $S_1$ at
the right end).

A structural proof would also explain the asymmetric clause in
($\gamma$): the non-continuation $9 \notin$ column $1$ should arise
as a transversality condition on the letter-shift dynamics at the
column-$1$ boundary.

\subsection{Correction to the memory note}
The memory note [[omega-survivor-pattern]] flagged a count of $12$
$\Om^*$-only zeros on $(3,3,1,1,1)$, characterised by row $1 = 1,2,3$
($11$ cases) or row $1$ starting $1, 2$ with row $2$ starting $3, 4$
adjacent ($1$ case). The direct probe finds $0$ $\Om^*$-only cases on
this shape: all $35$ residual zero diagonals satisfy $\Om v_T = 0$
\emph{and} $\Om^* v_T = 0$ simultaneously. The $12$ count is
retracted for $\la = (3,3,1,1,1)$. The original observation likely
referred to a different shape or to a finer right-kernel-only
stratum within the residual; this requires cross-checking against
the shape the memory was originally derived from.

The corrected accounting on $\la = (3,3,1,1,1)$ is:
\begin{center}
\begin{tabular}{l|c}
SR-at-$1$ residual ($p_T = 0$) & $35$ \\
of which $\Om v_T = 0$ & $35$ \\
of which $\Om^* v_T = 0$ & $35$ \\
of which $\Om^*$-only (i.e.\ $\Om v_T \ne 0$) & $0$ \\
\end{tabular}
\end{center}

\subsection{Why $\{3, 5, 7\}$?}
The set $\{3, 5, 7\}$ equals $\{q + 2k : k = 0, 1, 2\}$ with
$q = 3$. The chain factors $R'_3, R'_4, \ldots, R'_9$ index a
sequence of column-$1$ transition steps in the SRD reduction. The
killing stages $\{R'_6, R'_7, R'_8\}$ are the \emph{middle} of the
chain, not the endpoints; $R'_3$ (leftmost) and $R'_9$ (rightmost)
never appear as first-kill stages on the residual. This middle-of-
chain symmetry is suggestive but currently not exploited.

\section*{Acknowledgements}

This note builds on the SC-at-$1$ trivial reduction
(\texttt{2026-05-19-sc-at-1-trivial-reduction.tex}), the SR-conjecture
refutation (\texttt{2026-05-19-sr-conjecture-refuted.tex}), and the
two-sided kernel dichotomy
(\texttt{2026-05-18-d-class-row-zero.tex}). The Hoefsmit basis
machinery is at \texttt{2026-05-08-rank-Pi/verify\_seminormal.py};
diagonal probes at \texttt{2026-05-19-V33111-residual-diagnostic.py};
cell-partition data at
\texttt{2026-05-19-V332-V33111-cell-probe/probe33111.py}.

\begin{thebibliography}{9}

\bibitem{DaiZhang}
Dai, X.\ and Zhang, Z.,
\emph{Recursive Gelfand $W$-graphs and the canonical cells of the
symmetric group},
arXiv:2605.17844, 2026.

\bibitem{Hoefsmit}
Hoefsmit, P.\ N.,
\emph{Representations of Hecke algebras of finite groups with BN-pairs
of classical type},
Ph.D.\ thesis, University of British Columbia, 1974.

\bibitem{Murphy}
Murphy, G.\ E.,
\emph{On the representation theory of the symmetric groups and
associated Hecke algebras},
J.\ Algebra \textbf{152} (1992), 492--513.

\bibitem{LMRSW}
Liu, Y., Mazel-Gee, A., Reutter, D., Stroppel, C., and Wedrich, P.,
\emph{A braided $(\infty, 2)$-categorical structure on Soergel
bimodules},
arXiv:2401.02956, 2024.

\end{thebibliography}

\end{document}
